\subsection{Частные производные и дифференциалы высших порядков}

\begin{definition}
	\textit{Частная производная $k$-ого порядка для  данного упорядоченного набора переменных $x_{a_1},x_{a_2},...,x_{a_k}, a_i\in\{1,...,n\}$}, $k>1$ функции $f:\R^n\to\R$ определяется по индукции, как частная производная по переменной $x_{a_1}$, соответствующей производной $k-1$-ого порядка (при  этом cуществование $k$-ой производной вообще говоря не гарантируется):
	\[
		\frac{\vdelta^k f}{\vdelta x_{a_1} \vdelta x_{a_2}...\vdelta x_{a_k}} (\vv{x}_0) := \pd{}{x_{a_1}}\left(\frac{\vdelta^{k-1} f}{\vdelta x_{a_2} ...\vdelta x_{a_k}} \right)(\vv{x}_0)
	\]
	Стоит заметить, что $a_i$-ые совершенно не обязательно различны. Более того, если индексы нескольких подряд идущих переменных равны так что $a_m = a_{m+1}=...=a_{m+l} = b$, для удобства вместо $\vdelta x_{a_m} \vdelta x_{a_{m+1}}...\vdelta x_{a_{m+l}}$ пишут $\vdelta x_b^{l+1}$. Ну и, конечно для краткости записи берут наибольшие промежутки идущих подряд равных индексов. Кроме того, как мы увидим далее, порядок $\vdelta x_{a_i}$-ых вообще говоря важен, и переставлять их между собой просто так нельзя.\\
	\textbf{Пример.} Таким образом такую производную:
	\[
		\frac{\vdelta^4 f}{\vdelta x_{1} \vdelta x_{1}\vdelta x_{2}\vdelta x_{1}} (\vv{x}_0)
	\]
	наиболее правильно будет записать так:
	\[
	\frac{\vdelta^4 f}{\vdelta x_{1}^2 \vdelta x_{2}\vdelta x_{1}} (\vv{x}_0)
	\]
	\begin{anote}
		Напомним, что $\pd{}{x_j}$ - операторы. Аналогично можно смотреть на $	\frac{\vdelta^k }{\vdelta x_{a_1} \vdelta x_{a_2}...\vdelta x_{a_k}}$ как на соответствующий оператор, и опять же нашу запись можно просто считать синтаксическим сахаром: 
		\[
				\frac{\vdelta^k f}{\vdelta x_{a_1} \vdelta x_{a_2}...\vdelta x_{a_k}} = 	\frac{\vdelta^k }{\vdelta x_{a_1} \vdelta x_{a_2}...\vdelta x_{a_k}}(f)
		\]
		 Зададим на операторах операции сложения и умножения вот по таким правилам:
		\begin{itemize}
			\item Сложение: $(\alpha + \beta)(f) = \alpha(f) +\beta(f)$
			\item Умножение: $(\alpha  \beta)(f) = \alpha( \beta(f))$
			\item Умножение на функцию: $(\alpha g)(f) = \alpha(f)g$
		\end{itemize}
	то есть сложение - интуитивное, а умножение - просто композиция. Посмотрим, как запишется умножение в нынешней нотации операторов дифференцирования высших порядков:
	\[
	\pd{}{x_i}\pd{}{x_j}(f) := \pd{}{x_i}(\pd{}{x_j}(f)) =:\frac{\vdelta^2}{\vdelta x_i \vdelta x_j}(f)
	\]
	
	\begin{multline*}
		\frac{\vdelta^k }{\vdelta x_{a_1} \vdelta x_{a_2}...\vdelta x_{a_k}} 	\frac{\vdelta^l }{\vdelta x_{b_1} \vdelta x_{b_2}...\vdelta x_{b_l}}(f) :=
		\\
		\frac{\vdelta^k }{\vdelta x_{a_1} \vdelta x_{a_2}...\vdelta x_{a_k}} 	(\frac{\vdelta^l f}{\vdelta x_{b_1} \vdelta x_{b_2}...\vdelta x_{b_l}}(f))
		\\
		=:	\frac{\vdelta^{k+l} }{\vdelta x_{a_1} \vdelta x_{a_2}...\vdelta x_{a_k}\vdelta x_{b_1} \vdelta x_{b_2}...\vdelta x_{b_l}}(f)
	\end{multline*}
	Внезапно в нашей нотации умножение выглядит <<интуитивно>>, мы буквально получили в результате то же, что и если бы просто перемножили <<числители и знаменатели>> соответствующих формальных записей (в этом случае надо думать, что $\vdelta$ сверху и $\vdelta x_i$ снизу - переменные, при чём вторая <<неделимая>> и просто записывается в две буквы, прямо как дифференциал). По аналогичным причинам, естественно, работает запись подряд идущих одинаковых $\vdelta x_j$ как степень. Опять же неожиданный синтаксический сахар. Дальнейшее применение этого всего будет в формуле дифференциала высших порядков, там же и понадобится описанные выше сложение и умножение на функцию. \textbf{Конец замечания автора}
	\end{anote}
	Ещё одна форма обозначения таких производных такова:
	\[
		f_{x_{a_1},x_{a_2}...x_{a_k}}^{(k)}(\vv{x_0}):=\frac{\vdelta^k f}{\vdelta x_{a_1} \vdelta x_{a_2}...\vdelta x_{a_k}} (\vv{x}_0)
	\]
\end{definition}

\begin{example}
	Рассмотрим следующую функцию:
	\[
		f(x, y) = \System{
			&{xy \frac{x^2 - y^2}{x^2 + y^2},\ x^2 + y^2 \neq 0}
			\\
			&{0,\ x^2 + y^2 = 0}
		}
	\]
	Частная производная по $x$ имеет вид:
	\[
		\pd{f}{x} (x, y) = \System{
			&{y \frac{x^2 - y^2}{x^2 + y^2} + \frac{4x^2 y^3}{(x^2 + y^2)^2},\ x^2 + y^2 \neq 0}
			\\
			&{0,\ x^2 + y^2 = 0}
		}
	\]
	А по $y$ она выглядит так:
	\[
		\pd{f}{y} (x, y) = \System{
			&{x \frac{x^2 - y^2}{x^2 + y^2} - \frac{4x^3 y^2}{(x^2 + y^2)^2},\ x^2 + y^2 \neq 0}
			\\
			&{0,\ x^2 + y^2 = 0}
		}
	\]
	А теперь посчитаем вторые производные $f''_{xy}$ и $f''_{yx}$ в нуле:
	\[
		\frac{\vdelta^2}{\vdelta x \vdelta y}(0, 0) = \liml_{\Delta x \to 0} \frac{\pd{f}{y}(\Delta x, 0) - \pd{f}{y}(0, 0)}{\Delta x} = 1
	\]
	\[
		\frac{\vdelta^2}{\vdelta y \vdelta x}(0, 0) = \liml_{\Delta y \to 0} \frac{\pd{f}{x}(0, \Delta y) - \pd{f}{x}(0, 0)}{\Delta y} = -1
	\]
	Таким образом, значение производной зависит от порядка дифференцирования
\end{example}

\begin{theorem} (Шварц)
	Если существуют $\frac{\vdelta^2 f}{\vdelta x \vdelta y}$ и $\frac{\vdelta^2 f}{\vdelta y \vdelta x}$ в некоторой окрестности точки $(x_0, y_0)$, причём они непрерывны в $(x_0, y_0)$, то
	\[
		\frac{\vdelta^2 f}{\vdelta x \vdelta y} (x_0, y_0) = \frac{\vdelta^2 f}{\vdelta y \vdelta x} (x_0, y_0)
	\]
\end{theorem}

\begin{proof}
	Рассмотрим выражение следующего вида:
	\[
		\Phi (h) = f(x_0 + h, y_0 + h) - f(x_0 + h, y_0) - f(x_0, y_0 + h) + f(x_0, y_0)
	\]
	Чтобы преобразовать $\Phi(h)$, определим допольнительную функцию $\phi$:
	\[
		\phi(x) := f(x, y_0 + h) - f(x, y_0)
	\]
	Отсюда получаем новое выражение:
	\[
		\Phi(h) = \phi(x_0 + h) - \phi(x_0)
	\]
	Коль скоро $f$ дифференцируема в окрестности точки $(x_0, y_0)$, то и $\phi$ - тоже (при достаточно малых $h$). Более того, мы можем расписать функцию $\Phi(h)$, которая описывает приращение $\phi$, и после применить к $\phi$ теорему Лагранжа для $x \in (x_0; x_0 + h)$ (взяв достаточно малое $h>0$), найти $\xi\in (x_0; x_0 + h)$ такое что:
	\[
		\Phi(h) = \phi(x_0 + h) - \phi(x_0) = \phi'(\xi) h = \left(\pd{f}{x} (\xi, y_0 + h) - \pd{f}{x} (\xi, y_0)\right) h
	\]
	Теперь можно снова применить теорему Лагранжа к $\pd{f}{x} (\xi, y),\ y \in (y_0, y_0 + h)$, как функции от $y$ найти $\nu\in (y_0; y_0 + h)$ такое что:
	\[
		\left(\pd{f}{x} (\xi, y_0 + h) - \pd{f}{x} (\xi, y_0)\right) h = \frac{\vdelta^2 f}{\vdelta y \vdelta x} (\xi, \nu) h^2
	\]
	Теперь распишем функцию $\psi(y)$ следующего вида:
	\[
		\psi(y) := f(x_0 + h, y) - f(x_0, y)
	\]
	Отсюда снова получаем способ по-другому записать $\Phi(h)$, с аналогичными ограничениями на $\theta$ и $\eta$:
	\[
		\Phi(h) = \psi(y_0 + h) - \psi(y_0) = \psi'(\eta) h = \left(\pd{f}{y} (x_0 + h, \eta) - \pd{f}{y} (x_0, \eta)\right) h = \frac{\vdelta^2 f}{\vdelta x \vdelta y} (\theta, \eta) h^2
	\]
	Таким образом, при достаточно малых $h$ мы получили:
	\[
		\frac{\vdelta^2 f}{\vdelta y \vdelta x} (\xi, \nu) = \frac{\Phi(h)}{h^2}=\frac{\vdelta^2 f}{\vdelta x \vdelta y} (\theta, \eta)
	\]
	в силу ограничений на $\xi, \nu, \theta, \eta$ при устремлении $h$ к нулю, обе точки $(\xi, \nu),\ (\theta, \eta)$ стремятся к $(x_0, y_0)$, а в силу непрерывности обоих вторых частных производных в этой точке, устремляя $h$ к нулю, мы получаем, что они в ней равны. 
\end{proof}

\begin{definition}
	Функция $f$ называется \textit{дважды дифференцируемой} в точке $\vv{x}_0 \in \R^n$, если \underline{все} её частные производные дифференцируемы в $\vv{x}_0$.
\end{definition}

\begin{theorem} (Юнг)
	Если $f$ дважды дифференцируема в точке $(x_0, y_0)$, то
	\[
		\frac{\vdelta^2 f}{\vdelta x \vdelta y} (x_0, y_0) = \frac{\vdelta^2 f}{\vdelta y \vdelta x} (x_0, y_0)
	\]
\end{theorem}

\begin{proof}
	Положим $\Phi(h)$ и $\phi(x)$ как в предыдущем доказательстве. Тогда:
	\[
		\Phi(h) = \phi(x_0 + h) - \phi(x_0) = \phi'(\xi) h = \left(\pd{f}{x} (\xi, y_0 + h) - \pd{f}{x} (\xi, y_0)\right) h
	\]
	Прибавим и отнимем частную производную в точке $(x_0, y_0)$. Отсюда имеем:
	\[
		\Phi(h) = \left(\pd{f}{x} (\xi, y_0 + h) - \pd{f}{x} (x_0, y_0)\right) h - \left(\pd{f}{x} (\xi, y_0) - \pd{f}{x} (x_0, y_0)\right) h =
	\]
	Распишем по определению дифференцируемости полученные приращения:
	\begin{multline*}
		=\left(\frac{\vdelta^2 f}{\vdelta x^2} (x_0, y_0) (\xi - x_0) + \frac{\vdelta^2 f}{\vdelta y \vdelta x} (x_0, y_0) h + o\left(\sqrt{(\xi - x_0)^2 + h^2}\right)\right) h -
		\\
		-\left(\frac{\vdelta^2 f}{\vdelta x^2} (x_0, y_0) (\xi - x_0) + o(|\xi - x_0|)\right) h,\ h\to 0
	\end{multline*}
	Очевидно, что $|\xi - x_0| \leq |(\xi - x_0, h)| = \sqrt{(\xi - x_0)^2 + h^2} \leq \sqrt{2}|h|$.
	В итоге $\Phi(h)$ принимает вид
	\[
		\Phi(h) = \frac{\vdelta^2 f}{\vdelta y \vdelta x} (x_0, y_0) h^2 + o(h^2),\ h\to 0
	\]
	Точно так же рассматривается разложение $\Phi(h)$ через $\psi(y)$. Приравнивая эти разложения, получим
	\[
		\frac{\vdelta^2 f}{\vdelta y \vdelta x} (x_0, y_0) + \frac{o(h^2)}{h^2} = \frac{o(h^2)}{h^2} + \frac{\vdelta^2 f}{\vdelta x \vdelta y} (x_0, y_0)
	\]
	При стремлении $h \to 0$ это даёт нам требуемое равенство.
\end{proof}

\begin{example}
	Приведённые теоремы являются достаточными условиями, но не необходимыми. Данная функция иллюстрирует это:
	\[
		f_1(x, y) = \System{
			&{x^{3/2} y^{3/2} \sin \frac{1}{x} \sin \frac{1}{y},\ xy \neq 0}
			\\
			&{0,\ xy = 0}
		}
	\]
	Для определённости, будем рассматривать $x, y > 0$. Тогда, частная производная по $x$ имеет вид:
	\[
		\pd{f_1}{x} = \frac{3}{2} x^{1/2} y^{3/2} \sin \frac{1}{x} \sin \frac{1}{y} - \frac{y^{3/2}}{x^{1/2}} \cos \frac{1}{x} \sin \frac{1}{y},\ xy \neq 0
	\]
	В точках вида $(0, y)$ (то есть при $x = 0$) эта же производная будет иметь значение:
	\[
		\pd{f_1}{x} (0, y) = \liml_{\Delta x \to 0} \frac{f_1(\Delta x, y) - f_1(0, y)}{\Delta x} = 0
	\]
	Теперь посчитаем вторую частную производную по $y$ в ситуации $x, y > 0$:
	\[
		\frac{\vdelta^2 f_1}{\vdelta y \vdelta x} = \frac{9}{4} x^{1/2} y^{1/2} \sin \frac{1}{x} \sin \frac{1}{y} - \frac{3x^{1/2}}{2y^{1/2}} \sin \frac{1}{x} \cos \frac{1}{y} - \frac{3y^{1/2}}{2x^{1/2}} \cos \frac{1}{x} \sin \frac{1}{y} + \frac{1}{x^{1/2} y^{1/2}} \cos \frac{1}{x} \cos \frac{1}{y} 
	\]
	В нуле при этом смешанная производная будет равна
	\[
		\frac{\vdelta^2 f_1}{\vdelta y \vdelta x}(0, 0) = \liml_{\Delta y \to 0}\frac{\pd{f_1}{x}(0, \Delta y) - \pd{f_1}{x}(0, 0)}{\Delta y} = 0
	\]
	Аналогично можно показать, что 
	\[
		\frac{\vdelta^2 f_1}{\vdelta x \vdelta y}(0, 0) = \liml_{\Delta x \to 0}\frac{\pd{f_1}{y}(\Delta x, 0) - \pd{f_1}{y}(0, 0)}{\Delta x} = 0
	\]
	\begin{anote}
		\textit{(Для тех кто не понимает, почему посчитанные частные и смешанные производные не непрерывны в нуле)} Лукашов посчитал этот факт не требующим пояснения, но мы всё немного уточним.  
		\begin{itemize}
			\item Для $\pd{f_1}{x}$ можно рассмотреть последовательность $(x_n,y_n) = (\frac{1}{2\pi n^4},\frac{1}{2\pi n + \frac{\pi}{2}})$ несложно увидеть, что для неё первое слагаемое в выражении $\pd{f_1}{x}$ будет бесконечно малой, а второе неограниченно возрастать, и к нулю их разность не сойдётся.
			\item Для $
			\frac{\vdelta^2 f_1}{\vdelta y \vdelta x}$ можно рассмотреть последовательность $(x_n, y_n) = (\frac{1}{2\pi n}, \frac{1}{2\pi n})$ Тогда первые три слагаемых в выражении $
			\frac{\vdelta^2 f_1}{\vdelta y \vdelta x}$ будут ограниченны, а четвёртое - бесконечно большое. Значит опять же к нулю оно не сойдётся.
		\end{itemize}
	\end{anote}
	
	То есть $\frac{\vdelta^2 f_1}{\vdelta x \vdelta y} = \frac{\vdelta^2 f_1}{\vdelta y \vdelta x}$, хотя смешанные производные не непрерывны в нуле, частные производные тоже (а следовательно частные не дифференцируемы и сама функция не дважды дифференцируема). То есть обе теоремы являются лишь достаточными условиями.
\end{example}

\begin{example}
	Данная функция удовлетворяет теореме Юнга, но не удовлетворяет теореме Шварца
	\[
		f_2(x, y) = \System{
			&{x^3 y^3 \sin \frac{1}{xy},\ xy \neq 0}
			\\
			&{0,\ xy = 0}
		}
	\]
	
	Посчитаем частную производную по $x$:
	\[
		\pd{f_2}{x}(x, y) = \System{
			&{3x^2y^3 \sin \frac{1}{xy} - xy^2 \cos \frac{1}{xy},\ xy \neq 0}
			\\
			&{0,\ xy = 0}
		}
	\]
	Её дифференцируемость:
	\[
		\pd{f_2}{x}(x,y) \stackrel{?}{=} o(|(x,y)|),\ |(x,y)|\to0
	\]
	- верна, т.к. если расписать в полярных координатах, то оба слагаемых $\pd{f_2}{x}(x,y)$ будут содержать $r$ более чем первого порядка умноженную на ограниченные функции.
	\\\\
	Вторая частная производная по $y$:
	\[
		\frac{\vdelta^2 f_2}{\vdelta y \vdelta x}(x, y) = \System{
			&{9x^2y^2 \sin \frac{1}{xy} - 5xy\cos\frac{1}{xy} - \sin\frac{1}{xy},\ xy \neq 0}
			\\
			&{0,\ xy = 0}
		}
	\]
	Причем в нуле она имеет разрыв, так как $\sin\frac{1}{xy}$ не сходится.
\end{example}

\begin{example}
	Пример функции, которая удовлетворяет теореме Шварца, но не удовлетворяет теореме Юнга
	\[
		f_3(x, y) = \System{
			&{x^2 \sin \frac{1}{x},\ x \neq 0}
			\\
			&{0,\ x = 0}
		}
	\]
	Частная производная по $x$ имеет следующий вид:
	\[
		\pd{f_3}{x} = \System{
			&{2x \sin \frac{1}{x} - \cos\frac{1}{x},\ x \neq 0}
			\\
			&{0,\ x = 0}
		}
	\]
	Из-за $\cos \frac{1}{x}$ она не непрерывна и тем более не дифференцируема, однако вторая смешанная производная по $y$:
	\[
		\frac{\vdelta^2 f_3}{\vdelta y \vdelta x} = 0
	\]
	Так как функция не зависит от $y$. Аналогично получаем другую смешанную производную и следующее равенство:
	\[
		\frac{\vdelta^2 f_1}{\vdelta x \vdelta y} = 0 = \frac{\vdelta^2 f_1}{\vdelta y \vdelta x}
	\]
\end{example}

\begin{definition}
	Функция называется \textit{$k$ раз дифференцируемой} в точке $\vv{x}_0$ $k>1$, если её частные производные $(k - 1)$-го порядка дифференцируемы. (Заметим что согласуется с дважды дифференцируемостью.) Если $f$ дифференцируема $k$ раз в $\vv{x_0} \in R^n$, то её дифференциал $k$-го порядка определяется по индукции, как дифференциал от её дифференциала $k-1$-го порядка:
	\[
		d^k f(\vv{x}_0) := d(d^{k - 1} f)(\vv{x}_0)\ ,
	\] причём дифференциалы независимых переменных $dx_i, \  i =1,...,n$ - одни и те же при каждом взятии дифференциала. А во время самого взятия, на них нужно смотреть как на константы вплоть до получения ответа, и только в конце отождествить старые и новые $dx_i$-ые.
	
	\begin{anote}
		Читатель опять может заметить некоторую странность - как мы могли дифференцировать хотя бы два раза, если уже после первого функция перестала быть числовой, и стала бить в линейные отображения. В определении мы описали, что формульно мы называем дифференциалом высшего порядка, но как это честно определить? Для последнего стоит сначала изучить что такое дифференцирование в ЛНП (глава 6), и понять, что тогда корректно будет определён дифференциал $k$-го порядка в точке $x_0$, и он будет уже не линейной формой, а симметричной $k$-линейной формой, то есть симметричным полилинейным отображением $\R^n\times\dots\times\R^n\to\R$. Наконец, как симметричные билинейные формы имеют естественную биекцию с $n$-формами (то что квадратичные и билинейные имеют таковую изучалось в курсе алгебры и и геометрии). И вот эта $n$-форма, ака однородный многочлен $n$-ой степени от $dx_k$-ых и будет тем, что мы понимаем под $d^kf_{x_0}$. (Хотя правильнее было бы всё же говорить именно о полилинейной форме, но на практике именно $n$-формы интересны).
	\end{anote}
	
	\begin{anote}
		Читатель может заметить некоторую странность, дифференциал $k$-го порядка определён только для $k$ раз дифференцируемой функции, в то же время он требует наличия дифференциала $k-1$-го порядка, то есть, получается для корректности, наша функция должна была быть ещё и $k-1$ раз дифференцируема ? И так далее ? В сущности, для корректности нам достаточно показать, что из дифференцируемости $k$ раз следует дифференцируемость $k-1$ раз, $k>1$:
		\begin{proof}
			По определению, нам нужно вывести из дифференцируемости производных $k-1$-го порядка, то же для производных $k-2$-го порядка. Ну, и правда, раз производные $k-1$-го дифференцируемы, то они непрерывны, а значит, будучи непрерывными частными производными производных $k-2$-го порядка, они обеспечивают дифференцируемость последних.
		\end{proof}
	\end{anote}
	В случае, если функция зависит только от $n$ независимых переменных, то справедлива формула:
	\[
		d^k f(\vv{x}_0) = \left(\suml_{j = 1}^n \pd{}{x_j} dx_j\right)^k f(\vv{x}_0)
	\]
	\begin{proof}
		Докажем по индукции. База $k=1$ - просто по определению первого дифференциала. Переход:
		\begin{multline*}
			d^k f(\vv{x}_0)
			= 
			d(d^{k - 1} f)(\vv{x}_0) 
			= 
			d\bigl(\left(\suml_{j = 1}^n \pd{}{x_j} dx_j\right)^{k-1}f\bigr) (\vv{x}_0) 
			= 
			\left( \suml_{i = 1}^n \pd{}{x_i} \left(\suml_{j = 1}^n \pd{}{x_j} dx_j\right)^{k-1}f\cdot dx_i \right) (\vv{x_0})  
			=
			\\
			=\left(\suml_{i = 1}^n \left(\pd{}{x_i}dx_i\right) \left(\suml_{j = 1}^n \pd{}{x_j} dx_j\right)^{k-1}\right)f(\vv{x_0}) 
			=
			\left(\left(\suml_{i = 1}^n \pd{}{x_i}dx_i\right) \left(\suml_{j = 1}^n \pd{}{x_j} dx_j\right)^{k-1}\right)f(\vv{x_0}) 
			=
			\\
			=
			\left(\suml_{j = 1}^n \pd{}{x_j} dx_j\right)^k f(\vv{x}_0) 
		\end{multline*}
	\end{proof}
\end{definition}

\begin{anote}
	На самом деле доказательство факта выше было бы очевидно, если бы мы сразу обговорили кое-что. А именно, что на дифференциал тоже можно(и нужно) смотреть как на оператор (возвращающий функцию, которую мы называли первым дифференциалом). Во-первых, тогда в формуле $d^k f(\vv{x}_0) := d(d^{k - 1} f)(\vv{x}_0)$, <<$d^k$>> будет не формальным обозначением дифференциала $k$-го порядка, а просто будет читаться, как возведение оператора в степень. Во-вторых, если теперь вспомнить, что $d f(\vv{x}_0) = \left(\suml_{j = 1}^n \pd{}{x_j} dx_j\right) f(\vv{x}_0)$, а на языке операторов, $d = \suml_{j = 1}^n \pd{}{x_j} dx_j$, то искомая формула будет просто возведением этого равенства в степень! И никаких громоздких бессмысленных формул и индукций!  
\end{anote}

\begin{addition}
	Более того эту формулу можно раскрыть так:
	\[
		d^k f(\vv{x}_0) = \left(\suml_{j = 1}^n \pd{}{x_j} dx_j\right)^k f(\vv{x}_0) 
		= 
		\suml_{j_1 +...+j_n = k}\frac{\vdelta^k f}{\vdelta x_1^{j_1}...\vdelta x_n^{j_n}}(\vv{x_0})\frac{k!}{j_1!...j_n!}dx_1^{j_1}...dx_n^{j_n}
	\]
	Что следует, из того, как раскрывается бином (смотри курс ОКТЧ), если понять, почему в нашем случае значение высшей производной не зависит от порядка взятия составляющих её частных производных. Лукашов сказал, что это просто следует из  Теоремы Юнга. Автор считает необходимым на всякий случай написать следующее разъяснение.
	\begin{anote}
		Итак, нам необходимо показать, что если функция $k$ раз дифференцируема, то две частные производные высших порядков, с одинаковым набором переменных, по которым ведётся дифференцирование, но отличающееся непосредственно его порядком, равны. Для начала, из курса алгебры и геометрии мы знаем, что две перестановки одного и того же набора переменных можно перевести друг в друга, транспозициями соседних элементов. Поэтому нам достаточно доказать утверждение только для наборов отличающихся такой транспозицией. То есть что
		\[
			\frac{\vdelta^k f}{\vdelta x_{a_1}...\vdelta x_{a_{i-1}}\vdelta x_{a_i}\vdelta x_{a_{i+1}}\vdelta x_{a_{i+2}}...\vdelta x_{a_k}} (\vv{x}_0)
			=
			\frac{\vdelta^k f}{\vdelta x_{a_1}...\vdelta x_{a_{i-1}}\vdelta x_{a_{i+1}}\vdelta x_{a_{i}}\vdelta x_{a_{i+2}}...\vdelta x_{a_k}} (\vv{x}_0)
		\] 
		для этого достаточно проверить, что
		\[
			\frac{\vdelta^{k-i+1} f}{\vdelta x_{a_i}\vdelta x_{a_{i+1}}\vdelta x_{a_{i+2}}...\vdelta x_{a_k}} (\vv{x}_0)
			=
			\frac{\vdelta^{k-i+1} f}{\vdelta x_{a_{i+1}}\vdelta x_{a_{i}}\vdelta x_{a_{i+2}}...\vdelta x_{a_k}} (\vv{x}_0)
		\]
		И уже здесь, мы можем воспользоваться тем, что функция $k$-раз дифференцируемая, как следствие $k-i+1$ раз дифференцируемая, как следствие функция
		\[
			\frac{\vdelta^{k-i-1} f}{\vdelta x_{a_{i+2}}...\vdelta x_{a_k}} (\vv{x}_0)
		\]
		дважды дифференцируема, а значит искомое равенство её смешанных производных уже и правда верно по признаку Юнга.
	\end{anote} 
\end{addition}

\begin{note}
	\textit{(Отсутствие инвариантности формы дифференциалов высших порядков)} Если $x_1,..., x_n$ - функции независимых переменных, то
	\[
		d^2f=d\left(\suml_{j = 1}^n \pd{f}{x_j} dx_j\right) = \suml_{i=1}^n \suml_{j = 1}^n \frac{\vdelta^2 f}{\vdelta x_i \vdelta x_j} dx_idx_j + \suml_{j=1}^n \pd{f}{x_j}d^2x_j
	\]
\end{note}

\begin{note}
	Формула для дифференциала $k$-го порядка справедлива и в том случае, когда $x_j$-ые - линейные функции независимых переменных. (Тогда в формуле выше, как и в её "потомках" все $d^kx_j\ , k>1$ просто по определению выродятся в $0$, и останется <<инвариантная>> форма.)
\end{note}